\documentclass{bmsnkresume}
\usepackage{amsmath}

\name{SHASHANK BIRUR MANJUNATHA}
\professional{Embedded Software Engineer}
\contact{Email: \href{mailto:shashankbmanjunatha@gmail.com}{shashankbmanjunatha@gmail.com} | Phone: +49 17685648702 | Location: Stuttgart, DE}
\socials{GitHub: \href{https://github.com/shashaaankk}{/shashaaankk} | LinkedIn: \href{https://www.linkedin.com/in/shashank-birur-manjunatha-367061104}{/shashank-birur-manjunatha-367061104}}
\photo{potrait00.png}

\begin{document}

\newcommand{\columnsectiontitle}[1]{%
  {\cvemphsize\color{darktext}\MakeUppercase{#1}}\par
  \noindent{\color{lighttext}\rule{\linewidth}{0.2pt}}%
  \vspace{0.2em}%
}

% Define academicsitems environment
\newenvironment{academicsitems}{\begin{itemize}}{\end{itemize}}


\makeheader

\sectiontitle{Professional Experience}

\experience
  {Robert  Bosch  Gesellschaft  für  medizinische  Forschung  mbH}
  {BMT, University of Stuttgart, Germany}
  {Research Assistant - Segmentation}
  {March 2026 - Present}
\begin{itemize}
  \item Developing automated liver vessel segmentation from contrast CT for interventional planning, combining pretrained foundation models with topology-aware post-processing
  \item Fine-tuning a vessel segmentation foundation model (vesselFM) on public liver CT datasets
\end{itemize}
\vspace{-0.6em}

\experience
  {Valeo Schalter und Sensoren GmbH}
  {Bietigheim-Bissingen, Germany}
  {Working Student - Driver Intention Determination}
  {Jul 2024 - Dec 2024}
\begin{itemize}
  \item Studied driver acceptance of L2 lane keeping. Identified trust and takeover delay as common pain points 
  \item Traced pain points to plausible scenarios in the requirements. Designed test scenarios, data acquisition strategies and developed Qt-tools for driver studies on test tracks and simulators
  \item Processed raw CAN traces from drives into training data and built a Hidden Markov Model that classifies lane departures as intended or unintended from vehicle dynamics signals alone at 68\% precision, without driver monitoring hardware
\end{itemize}
\vspace{-0.6em}

\experience
  {Continental Autonomous Mobility (India) Pvt. Ltd.}
  {Bengaluru, India}
  {Engineer - Vehicle Dynamics}
  {Oct 2021 - Sep 2023}
\begin{itemize}
  \item Developed and shipped C++ algorithms for steering angle offset, lateral acceleration estimation, and velocity correction
  \item Delivered ASPICE SWE.3/SWE.4 compliance for the module, including detailed design, state behavior documentation, and unit tests, achieving 95\%+ branch and MC/DC coverage
  \item Validated software via SIL benchmarking and on-target processor-in-the-loop testing with CANoe and Trace32; adapted component interfaces for backward compatibility with legacy recordings
  \item Prototyped radar-based lateral velocity correction using forward Doppler radar and IMU motion compensation
  % \item Performed root-cause analysis of failed black-box test scenarios and automated their acceptance tests
\end{itemize}
\vspace{-0.6em}

  \sectiontitle{Education}
  \educationentry
    {M.Sc. Information Technology (Embedded Systems Engineering)}
    {University of Stuttgart, Germany}
    {Oct 2023 - Sep 2026}
    {GPA: 1.8}
    {}

  \educationentry
    {Certificate of Advanced Studies in Artificial Intelligence}
    {Artificial Intelligence Software Academy, University of Stuttgart}
    {Feb 2025 - Aug 2025}
    {}
    {}

  \educationentry
    {B.E Electronics \& Communication Engineering}
    {P.E.S College of Engineering (Affiliated to VTU), Mandya}
    {Aug 2017 - July 2021}
    {GPA: 1.5}
    {}

\vspace{-0.5em}
\sectiontitle{ Skills}
\vspace{0.2em}

{\color{darktext} Technical: C/C++, Python, MATLAB/Simulink, Qt (QML), FreeRTOS, ROS2, Linux, CAN, TensorRT, NVIDIA Jetson\par}
{\color{darktext} Tools \& Process: CMake, GTest, PyTest, Docker, GitHub Actions, CANoe, Trace32, IBM DOORS, Jira, Confluence, ASPICE, Scrum\par}
{\color{darktext} Languages: English (C1), German (A2), Kannada (Native), Hindi (Fluent)\par}

\sectiontitle{Master Thesis}
\academicsentry
  {Risk-Aware Autonomy: An Edge Framework for Operational Risk Assessment}
  {Nov 2025 - May 2026}
  {}
  {Institute of Automation Technology and Software Systems, University of Stuttgart}
\begin{academicsitems}
  \item Built scene graphs with YOLO-E and BoT-SORT on an NVIDIA Jetson over ROS2 for open-world navigation planning
  \item Implemented a hybrid RAG pipeline with Neo4j KG and vector databases to assess operational risk at the edge
  \item Benchmarked LLM throughput across scales and architectures, backends (TensorRT-LLM, vLLM, and llama.cpp) on Jetson AGX Orin and Thor 
  \item Identified memory bandwidth bottlenecks to be the limiting factors for real-time Operational Risk Assessment
\end{academicsitems}

\newpage

\sectiontitle{Projects}
\vspace{0.65em}
\projectentry
  {Ball-Balancing Robot to Demonstrate the Effects of 5G Delay in Networked Control Systems}
  {C++, ESP32, FreeRTOS, UDP, Python}
\begin{itemize}
  \item Built a distributed robot control system on ESP32/FreeRTOS nodes with inter-process and network communication over UDP sockets between sensor, compute, and actuator nodes on a custom parallel manipulator
  \item Developed system-level C++ firmware for the sensing and actuation control loops, scheduled at 2 ms and 4 ms periods
  \item Measured end-to-end latency via packet timestamping under emulated packet loss and jitter; control held through 10 ms of added delay and broke down by 20 ms
\end{itemize}
\vspace{0.37em}

\projectentry
  {Autonomous Chess-Playing Robot Arms in Simulation}
  {Python, MuJoCo, PyTorch, NVIDIA NIM (Llama 3.2), NumPy}
\begin{itemize}
  \item Built a perceive-decide-act loop in which two ViperX 6-DoF arms play chess against each other in MuJoCo, from overhead camera perception to pick-and-place manipulation
  \item Trained a 13-class CNN piece classifier on procedurally generated single-piece renders to 99.6\% test accuracy, then fine-tuned it on crowded scenes with exact labels by construction, which removed full-board misreads
  \item Derived closed-form inverse kinematics and checked it against MuJoCo forward kinematics, and made grasping reliable through convex mesh decomposition with CoACD, measured gripper jaw geometry, and tuned contact damping
  \item Integrated an LLM move engine through NVIDIA NIM by serializing perception output to a board-state JSON, prompting for UCI moves, and running turn-taking under semaphores with full game logging
\end{itemize}
\vspace{0.37em}

\projectentry
  {Attitude Estimation on ARM Cortex-M4}
  {Embedded C++, nRF52840, PlatformIO, I2C, Python}
\begin{itemize}
  \item Built modular firmware to stream 9-DoF IMU data over serial, running a complementary filter, Madgwick, and an implemented Kalman filter with gyro-bias estimation over identical sensor samples
  \item Characterized angle random walk, bias instability, and rate random walk sensor noise via overlapping Allan deviation and applied hard and soft iron magnetometer calibration
  \item Developed host-side tooling for reproducible data capture, live signal visualization, and side by side logging of the estimator outputs
\end{itemize}

\sectiontitle{Additional Courses}
\vspace{0.65em}
{\color{darktext}Introduction to Embedded Machine Learning, Edge Impulse (Coursera), Oct 2025\par}
{\color{darktext}Optimizing Generative AI on Arm Processors, Arm (Coursera), Oct 2025\par}

\sectiontitle{Publication}
\vspace{0.65em}
{\color{darktext}Seed Quantifying and Packaging Unit for Small-Scale Enterprises. In \textit{Futuristic Trends in Electronics \& Instrumentation Engineering}, IIP Series, March 2024. DOI: \href{https://doi.org/10.58532/V3BIEI3P7CH1}{10.58532/V3BIEI3P7CH1}\par}

\end{document}
